\documentclass[journal]{IEEEtran}

\usepackage{amsmath,amssymb,amsfonts}
\usepackage{algorithmic}
\usepackage{algorithm}
\usepackage{graphicx}
\usepackage{textcomp}
\usepackage{xcolor}
\usepackage{booktabs}
\usepackage{multirow}
\usepackage{url}
\usepackage{microtype}
\usepackage{cite}

\begin{document}

\title{AgriLink: An Empirical, Logistics-Aware Decision Framework for Agricultural Wholesale Market Selection}

\author{Pechetti~Mounika,
        P.~Kiran~Kumar,~\IEEEmembership{Student Member,~IEEE},
        B.~Bhargav,
        S.~RavindraNath,
        and~S.~Kranthi~Kowshik%
\thanks{Prof. P. Mounika is with the School of Computer Science and Engineering (SCOPE), VIT-AP University, Amaravati, Andhra Pradesh 522237, India (e-mail: pechetti.mounika@vitap.ac.in). She served as the faculty project guide and research supervisor.}%
\thanks{P. Kiran Kumar (Reg. No. 23BCE9610), B. Bhargav (Reg. No. 23BCE9623), S. RavindraNath (Reg. No. 23BCE9516), and S. Kranthi Kowshik (Reg. No. 23BCE9611) are undergraduate student researchers with the School of Computer Science and Engineering (SCOPE), VIT-AP University, Amaravati, Andhra Pradesh 522237, India (e-mail: pasupulakiran250@gmail.com).}%
\thanks{The authors confirm that empirical AGMARKNET datasets, simulation code, and anonymized primary field survey data supporting the findings of this study are available upon reasonable request.}%
}

\markboth{IEEE Access, Preprint Draft, November 2026}%
{Mounika \MakeLowercase{\textit{et al.}}: AgriLink: Logistics-Aware Wholesale Market Selection}

\maketitle

\begin{abstract}
Smallholder farmers in developing agrarian economies face an acute operational dilemma when commercializing harvest surpluses: selling at the geographically nearest wholesale assembly yard (\textit{mandi}) minimizes road haulage mileage but frequently locks producers into localized harvest-time price slumps, while chasing headline prices broadcast by digital market information systems ignores non-linear rural transport tariffs and risks negative spatial arbitrage. In this paper, we evaluate \textbf{AgriLink}, an empirical, logistics-aware market selection framework that couples retrospective wholesale price distributions with road distance modeling, empirical highway circuity corrections ($C_f = 1.28$), and commercial freight tariff functions directly into the farm-gate dispatch decision rule. Parameterized by 2,304 validated AGMARKNET market-day transaction records across 18 regulated APMC wholesale markets in Andhra Pradesh over a 364-day active trading evaluation period (August 16, 2024 to August 14, 2025; excluding the August 15 Independence Day holiday closure) spanning five commodity archetypes (Maize, Groundnut, Mango, Cotton, and Arhar), our primary evaluation analyzes a pooled multi-seed cohort of 5,000 synthetic dispatch scenarios across five independent random spatial seeds. Across the pooled cohort, AgriLink achieves a mean net realization of \textbf{₹153,868.54}, delivering a statistically significant gain of \textbf{+₹13,263.18 (+9.43\%, pooled 95\% CI: [+₹12,634.80, +₹13,887.82]; cluster mean difference +₹10,751.27, cluster $t = 5.62, p < 0.001$, cluster 95\% CI: [+₹6,540.67, +₹14,961.87])} over the status-quo Nearest Market heuristic (S1) and \textbf{+₹1,573.23 (+1.03\%, pooled 95\% CI: [+₹1,380.30, +₹1,752.81]; cluster mean difference +₹1,694.93, $t = 5.30, p < 0.001$, cluster 95\% CI: [+₹991.05, +₹2,398.81])} over the Naive Highest Price heuristic (S2), capturing 95.56\% of the theoretical Ex-Post Oracle arbitrage frontier with an average absolute regret of ₹615.68 (0.40\%). In divergent scenarios where AgriLink selects a more profitable non-local market, it saves an average of 91.11~km of one-way transit mileage relative to the Naive Highest Price heuristic (S2) across all dispatches and 372.78~km across divergent scenarios (while incurring an incremental 89.04~km relative to Nearest Market to capture profitable spatial arbitrage). The framework's superiority remains strictly positive under one baseline linear freight model and four alternative, independent freight-cost models (slab tariffs, trip minimums, stochastic freight shocks, and empty-backhaul penalties) as well as institutional mandi friction deductions. Primary commercial freight quotes ($n = 25$) and prospective observational farmer dispatch tracking ($n = 60$) corroborate the empirical calibration.
\end{abstract}

\begin{IEEEkeywords}
Agricultural logistics, AGMARKNET, APMC mandis, circuity factor, cluster-robust inference, decision support systems, Light Commercial Vehicles (LCV), out-of-sample validation, spatial arbitrage, wholesale price dispersion.
\end{IEEEkeywords}

\section{Introduction}
\IEEEPARstart{P}{rimary} wholesale agricultural commercialization in developing agrarian economies is mediated by regulated physical assembly markets known as Agricultural Produce Market Committees (APMCs) or \textit{mandis}~\cite{sekhar2021agricultural}. Over the past two decades, public agencies and agricultural ministries have invested heavily in digital market information systems (MIS), most notably the electronic National Agriculture Market (e-NAM) and the national AGMARKNET portal~\cite{sekhar2021agricultural, agmarknet2024}. The foundational economic thesis underpinning these digital price dissemination platforms is the spatial \textit{Law of One Price} (LOP)~\cite{takayama1971spatial, badiane1998spatial}. Classical spatial trade theory posits that broadcasting real-time wholesale price intelligence lowers information search frictions, eliminates spatial price dispersion across regional markets, and empowers agricultural producers to capture spatial arbitrage premiums~\cite{takayama1971spatial}.

However, empirical price monitoring across developing rural assembly yards demonstrates that spatial price dispersion remains substantial, structural, and persistent~\cite{mitra2018asymmetric, chatterjee2023market}. In our validated AGMARKNET wholesale dataset across Andhra Pradesh mandis, modal clearing prices on identical trading dates for standard-grade agricultural commodities within an operational 150~km haulage radius diverge by 12.4\% to 38.6\% (measured as the mean coefficient of variation and min-max spread across candidate yards). Smallholder producers---who operate 86.2\% of Indian agricultural holdings~\cite{agricensus2018}---struggle acutely to capitalize on these price spreads. At harvest, producers face severe liquidity constraints, lack access to on-farm temperature-controlled storage, and must negotiate ad-hoc road haulage in fragmented, informal spot freight markets~\cite{chatterjee2023market, casaburi2013rural}.

Under these operational constraints, smallholders overwhelmingly default to one of two intuitive heuristic decision rules~\cite{goyal2010information, fafchamps2012impact}:
\begin{enumerate}
    \item \textit{The Nearest-Market Heuristic} (S1, $\arg\min D$): Constrained by perceived transport barriers, limited vehicle access, and relational credit ties with local village middlemen (\textit{arthiyas}), producers deliver crops to the geographically closest yard~\cite{mitra2018asymmetric}. While this eliminates transit mileage and logistics uncertainty, it locks producers into localized price slumps during peak harvest flushes~\cite{chatterjee2023market}.
    \item \textit{The Naive Highest-Price Heuristic} (S2, $\arg\max P$): Informed by mobile price bulletins, SMS broadcasts, or peer networks, farmers direct their crop toward whichever market reports the highest headline price~\cite{goyal2010information}. However, because point-to-point spot haulage by rural Light Commercial Vehicles (LCVs) scales steeply with distance, the incremental transit cost frequently outstrips the gross price advantage, resulting in \textit{negative spatial arbitrage}~\cite{chatterjee2023market, takayama1971spatial}.
\end{enumerate}

\subsection{Research Problem Statement}
The core research problem addressed in this study is the \textit{logistics-agnostic failure of digital agricultural market recommendations}. While existing digital public infrastructure provides agricultural producers with real-time price signals, it provides zero guidance regarding whether spatial price differentials exceed the physical cost of vehicle charter, distance-decay spoilage, and institutional mandi transaction fees. Consequently, smallholder farmers systematically suffer from either spatial lock-in (under-switching) or catastrophic freight expenditure (over-switching).

\subsection{Research Objectives}
To resolve this empirical and algorithmic problem, this study pursues four specific objectives:
\begin{enumerate}
    \item Formulate a computationally efficient, logistics-aware market selection decision rule coupling retrospective price distributions with circuity-corrected road freight models.
    \item Implement a leakage-free, prospective forward evaluation framework across 5,000 synthetic dispatch scenarios spanning five commodity archetypes in Andhra Pradesh.
    \item Validate the stability of algorithmic recommendations against one baseline and four alternative, non-linear freight models to prevent circularity on model assumptions.
    \item Ground theoretical parameters through primary commercial freight quotes ($n = 25$) and prospective observational farmer dispatch tracking ($n = 60$).
\end{enumerate}

\subsection{Research Questions and Testable Hypotheses}
We formulate four formal research questions (RQs) and three directional hypotheses (Hs):
\begin{itemize}
    \item \textbf{RQ1 (Spatial Arbitrage):} Does logistics-aware market selection achieve statistically significant improvements in net farm-gate realization compared to nearest-market dispatch?
    \item \textbf{RQ2 (Naive Price Trap):} Does AgriLink prevent the negative spatial arbitrage incurred by unhedged highest-price dispatch while reducing unnecessary road transit?
    \item \textbf{RQ3 (Freight Model Independence):} Can the proposed framework maintain its economic advantage when evaluated under non-linear, stochastic, and backhaul-penalized freight tariffs?
    \item \textbf{RQ4 (Commodity Heterogeneity):} How does crop value density govern spatial switching frequency and net revenue gains?
\end{itemize}

\textbf{Hypothesis 1 ($H_1$):} Logistics-aware market selection yields a statistically significant positive net realization premium over the nearest-market heuristic ($\mu_{\text{AgriLink}} - \mu_{\text{Nearest}} > 0$).

\textbf{Hypothesis 2 ($H_2$):} Naive highest-price dispatch produces negative spatial arbitrage in over 15\% of dispatches, whereas AgriLink curtails loss frequency to below 5\%.

\textbf{Hypothesis 3 ($H_3$):} The net realization gains of AgriLink remain statistically significant ($p < 0.05$) under independent slab, minimum-charge, and stochastic freight models.

\subsection{Methodological Deficiencies in Extant Literature}
Prior computational and economic literature suffers from four major limitations:
\begin{enumerate}
    \item \textit{Logistics Agnosticism}: Most computational agricultural DSS treat market selection as a time-series price forecasting problem evaluated via statistical metrics (e.g., RMSE), ignoring transport costs entirely.
    \item \textit{Circular Evaluation}: Extant logistics-aware models optimize decisions using a linear freight formula and evaluate outcomes using that exact same linear formula, creating a tautological advantage.
    \item \textit{Retrospective Data Leakage}: Studies routinely employ overlapping time slices or condition test samples on markets known to trade in the future (\textit{common cohort filtering}), violating temporal causality~\cite{kaufman2012leakage}.
    \item \textit{Lack of Grounded Primary Data}: Prior work relies almost exclusively on unverified parameter assumptions without primary field transporter rate surveys.
\end{enumerate}

\subsection{Core Contributions}
To the best of our knowledge, this study presents an empirical decision-support framework that integrates retrospective mandi price distributions, road-distance correction, commercial freight estimates, and institutional deductions into a logistics-aware wholesale market selection workflow for the evaluated Andhra Pradesh APMC network. This paper delivers six core contributions:
\begin{enumerate}
    \item \textbf{Empirical Decision Formulation}: An analytical dispatch rule coupling retrospective price distributions with regional road circuity ($C_f = 1.28$) and commercial haulage tariffs.
    \item \textbf{Leakage-Free Forward Protocol}: Evaluation across 5,000 synthetic scenarios using strictly prior 30-day lookback series evaluated against forward auction prices over a 364-day active trading period.
    \item \textbf{Independent Freight Scoring}: Definitive elimination of circularity by evaluating performance under one baseline linear model and four alternative external freight models.
    \item \textbf{Institutional Frictions \& Perishable Spoilage}: Comprehensive ex-post modeling incorporating statutory APMC cess, commission, hamali handling, tare, queue penalties, and commodity-specific transit spoilage.
    \item \textbf{Cluster-Robust Statistical Inference}: Robust two-way clustered $t$-tests and block bootstrap confidence intervals across 12 monthly temporal blocks ($G = 12, \text{df} = 11$), reporting both pooled and cluster-based confidence intervals.
    \item \textbf{Primary Empirical Grounding}: Field grounding via 25 commercial transporter quotes, 60 observational farmer tracking cases, and 7 FPO pilot deployments (32 collective shipments).
\end{enumerate}

\section{Institutional Setting and Related Work}
\subsection{The Mandi Wholesale System in Andhra Pradesh}
Agricultural trade in Andhra Pradesh is regulated under the Andhra Pradesh (Agricultural Produce and Livestock) Markets Act. Wholesale primary transactions occur across regulated physical assembly yards (\textit{APMC mandis})~\cite{sekhar2021agricultural}. In Andhra Pradesh, primary transactions incur statutory deductions: Market Cess (1.0\%), Commission/Arthiya charges (1.0\%), and Hamali/weighing charges (₹0.20/kg). While the e-NAM platform enables electronic bidding, physical transport remains decentralized~\cite{sekhar2021agricultural}. Farmers charter Light Commercial Vehicles (LCVs, 1.5 to 3.5-tonne payload) on informal spot markets where transporters charge round-trip rates to cover empty backhauls.

\subsection{Spatial Price Dispersion and Information Frictions}
Neoclassical spatial equilibrium theory posits that under zero trade costs and perfect information, spatial price differentials cannot exceed physical transport costs ($|P_i - P_j| \le T_{ij}$)~\cite{takayama1971spatial, fackler2001spatial}. However, randomized evaluations demonstrate that price information alone does not automatically translate into higher farmer revenues. In a randomized controlled trial of SMS price alerts across 100 villages in Maharashtra, Fafchamps and Minten (2012)~\cite{fafchamps2012impact} observed that providing market price broadcasts did not result in statistically significant increases in prices received by farmers relative to control villages. Similarly, Mitra et al. (2018)~\cite{mitra2018asymmetric} found that providing potato farmers in West Bengal with periodic SMS price information did not alter average farm-gate sales prices or middleman margins. Chatterjee (2023)~\cite{chatterjee2023market} demonstrated that spatial price dispersion across Indian mandis remains wide and persistent, with intermediaries extracting margins when transport barriers restrict smallholder geographic mobility. These empirical findings motivate our framework: simply broadcasting prices does not empower farmers unless paired with decision support that accounts for the physical costs of transportation.

\subsection{Rural Haulage Economics and Road Circuity}
Geodesic straight-line distance ($D_g$) systematically underestimates road travel distance due to topography and road curvature~\cite{ballou2002selected, boscoe2012nationwide}. We relate road transit distance $D_r$ to straight-line distance $D_g$ via an empirical circuity factor $C_f$:
\begin{equation}
D_r = C_f \cdot D_g
\label{eq:circuity}
\end{equation}
Empirical route analysis across 153 inter-mandi routes connecting all 18 APMC mandis in Andhra Pradesh (calculated by comparing Haversine distance with actual road routing from OpenStreetMap/OSRM highway networks) reveals a mean circuity ratio of $C_f = 1.28 \pm 0.05$ (median: 1.27, range: 1.14 to 1.42), which we adopt as the regional baseline.

\begin{table*}[t]
\caption{Comparative Analysis of Agricultural Market Information and Decision-Support Systems}
\label{tab:lit}
\centering
\begin{tabular}{lcccccc}
\toprule
\textbf{Study / System} & \textbf{Year} & \textbf{Platform / Focus} & \textbf{Freight-Aware} & \textbf{Real-Time Dispatch} & \textbf{Spatial Switching} & \textbf{Primary Field Grounding} \\
\midrule
AGMARKNET Portal~\cite{agmarknet2024} & 2000--24 & Institutional reporting portal & No & No & No & No (Secondary registry data) \\
e-NAM Platform~\cite{sekhar2021agricultural} & 2016--24 & Electronic trading platform & No & No & Partial & No (Yard transactions) \\
Goyal~\cite{goyal2010information} & 2010 & Kiosk price access study & No & No & Village-level & Quasi-experimental survey \\
Fafchamps \& Minten~\cite{fafchamps2012impact} & 2012 & SMS price broadcast RCT & No & No & No & RCT (100 villages, Maharashtra) \\
Mitra et al.~\cite{mitra2018asymmetric} & 2018 & SMS alert field trial & No & No & No & RCT (72 villages, West Bengal) \\
Negi et al.~\cite{negi2018farmers} & 2018 & Channel choice econometric study & Partial & No & No & Cross-sectional survey (4 states) \\
Chatterjee~\cite{chatterjee2023market} & 2023 & Spatial equilibrium model & Modeled & No & Structural model & Secondary mandi database \\
\textbf{AgriLink (This Study)} & \textbf{2026} & \textbf{Logistics-aware dispatch framework} & \textbf{Yes (5 Models)} & \textbf{Yes (Algorithmic)} & \textbf{Yes (18 Mandis)} & \textbf{Yes (25 quotes, 60 cases)} \\
\bottomrule
\end{tabular}
\end{table*}

\subsection{Identified Research Gap}
Existing market information systems focus exclusively on unilateral price broadcasts or retrospective econometric estimation. No prior operational framework couples dynamic retrospective price distributions with circuity-corrected transport freight tariffs and institutional friction modeling to optimize farm-gate dispatch decisions.

\section{Mathematical Formulation and System Architecture}
\subsection{System Architecture and Processing Pipeline}
The AgriLink system architecture comprises six operational modules: 1) Farmer Input Module captures coordinates $L_f$, commodity $c$, parcel $Q \in [1,000, 3,500]$~kg, and date $T$; 2) Candidate Mandi Retrieval queries mandis within a geodesic screening radius $D_{\text{g,max}} = 150$~km; 3) Circuity-Corrected Distance Engine computes road distance $D_r$; 4) Price Estimator converts AGMARKNET prices from ₹/quintal to ₹/kg and computes historical reference price over lookback $[T-30, T-1]$; 5) Freight Engine evaluates commercial line-haul costs and institutional deductions; 6) Recommendation Engine ranks candidate mandis by expected net realization.

\subsection{Price Unit Conversion}
AGMARKNET administrative records report prices in ₹/quintal (where 1 agricultural quintal $\equiv$ 100 kg). Prior to algorithmic processing, all modal wholesale prices are converted to ₹/kg:
\begin{equation}
P_{\text{₹/kg}} = \frac{P_{\text{₹/quintal}}}{100}
\label{eq:price_conv}
\end{equation}
All subsequent equations, algorithms, optimization routines, tables, and reported realization values throughout this paper operate strictly in ₹/kg.

\subsection{Decision Timestamp \& Leakage Elimination}
In physical APMC wholesale markets, auctions commence at 09:30 AM. Farm-gate loading and vehicle dispatch decisions must be formed at \textbf{06:00 AM on auction date $T$}. At this timestamp, same-day auction outcomes are physically non-existent. To strictly enforce temporal causality and eliminate leakage~\cite{kaufman2012leakage}, historical clearing prices are sampled strictly within:
\begin{equation}
W_T = [T - 30, T - 1]
\label{eq:lookback}
\end{equation}

\subsection{Geodesic Screening vs. Road Haulage Corridor}
Candidate mandis are retrieved using a 150~km geodesic radius ($D_g \le 150$~km) measured from the farm-gate location $L_f$. Applying the regional circuity factor ($C_f = 1.28$), this geodesic screening boundary corresponds to an operational road corridor of up to $D_{\text{r,max}} = C_f \cdot D_{\text{g,max}} = 1.28 \times 150 = \mathbf{192\text{ km}}$. Candidate filtering occurs on $D_g$ to rapidly prune candidate search space, while all haulage cost calculations operate strictly on circuity-corrected road distance $D_r$.

\subsection{Two-Stage Mathematical Formulation}
We clearly delineate between the real-time Market Selection stage (Stage A) and the post-dispatch Final Net Realization evaluation stage (Stage B):

\textit{Stage A (Market Selection Optimization):} Let $L_f = (\phi_f, \lambda_f)$ denote farmer coordinates and $Q$ denote parcel weight in kg. Geodesic distance to market $m$ is:
\begin{equation}
\Delta\phi = \phi_m - \phi_f, \quad \Delta\lambda = \lambda_m - \lambda_f
\end{equation}
\begin{equation}
a = \sin^2(\Delta\phi/2) + \cos(\phi_f)\cos(\phi_m)\sin^2(\Delta\lambda/2)
\end{equation}
\begin{equation}
D_g = 2R \cdot \arcsin(\sqrt{a}), \quad D_r = C_f \cdot D_g
\end{equation}
where $R = 6,371$~km and $C_f = 1.28$. The historical reference price is:
\begin{equation}
\bar{P}_{m,c}(T) = \frac{1}{|T_{m,c}|} \sum_{t \in T_{m,c}} P_{m,c}(t), \quad T_{m,c} \subseteq [T-30, T-1]
\end{equation}
Baseline haulage cost is parameterized as:
\begin{equation}
TC(Q, D_r) = Q \cdot D_r \cdot c_{\text{rate}}
\end{equation}
where $c_{\text{rate}} = \text{₹}0.010/\text{kg}\cdot\text{km}$ (₹10.00/t-km). AgriLink selects target market $m^*$ optimizing expected Net Realization:
\begin{equation}
m^* = \arg\max_{m \in M} \left( Q \cdot \bar{P}_{m,c}(T) - TC(Q, D_r) \right)
\end{equation}

\textit{Stage B (Final Realized Net Earnings Evaluation):} Out-of-sample forward evaluation is conducted on observed forward prices at horizon $H = 7$ days (with sensitivity tested at $H \in \{14, 30\}$ days):
\begin{equation}
NR_{\text{realized}}(m^*) = Q \cdot P_{m^*,c}(T+H) - TC_{\text{eval}}(Q, D_r)
\end{equation}
Statutory yard deductions include Cess (1.0\%), Commission (1.0\%), Hamali (₹0.20/kg), Tare (1.5\%), and Queue Delay (₹350 for $D_r > 100$~km). Perishable Mango spoilage decay is:
\begin{equation}
L_{\text{spoil}}(D_r) = \min\left(0.12, \frac{D_r}{50} \times 0.0075\right) \cdot \text{Gross Revenue}
\end{equation}
\begin{equation}
\Pi_{\text{net}} = NR_{\text{realized}} - \sum \text{Frictions}
\end{equation}

\textit{Methodological Note on Objective Separation:} Stage A optimizes expected price less freight because only price histories and transport tariffs are observable to the farmer at 06:00 AM prior to departure. Mandi-specific unloading delays, local trader grading deductions, and variable queue penalties are ex-post realizations that cannot be reliably forecast ex-ante. Consequently, Stage B subjects the chosen market to comprehensive real-world institutional deductions and alternative freight tariff structures. We acknowledge that because statutory cess and queue delays vary across yards, the market chosen in Stage A may not invariably maximize the ex-post realization in every realized shipment---a realistic reflection of decentralized market friction.

\subsection{External Freight Scoring Models}
To eliminate circularity, realized outcomes are evaluated under one baseline model (Model A) and four alternative freight models:
\begin{enumerate}
    \item \textbf{Model A (Linear Baseline):} $TC = Q \cdot D_r \cdot 0.010$, reflecting standard volume chartering.
    \item \textbf{Model B (Distance Slab Tariff):} ₹0.0150/kg$\cdot$km for $D_r \le 50$~km, ₹0.0100/kg$\cdot$km for $50 < D_r \le 120$~km, and ₹0.0080/kg$\cdot$km for $D_r > 120$~km, reflecting long-haul economies of scale.
    \item \textbf{Model C (Minimum Trip Charter):} $\max(\text{₹}1,500, Q \cdot D_r \cdot 0.010)$ applied per trip, capturing minimum transporter call-out fees.
    \item \textbf{Model D (Stochastic Freight Shocks):} A random rate $c \sim U(0.010, 0.030)$ sampled \textbf{per dispatch trip}, simulating diesel surges and spot market tightening.
    \item \textbf{Model E (Backhaul Penalty):} $+50\%$ freight surcharge applied when $D_r > 75$~km, reflecting the 75~km intra-district boundary beyond which return cargo is uncertain and drivers charge deadhead return mileage.
\end{enumerate}

\subsection{Statistical Inference \& Dual Estimand Reporting}
Because wholesale mandi prices exhibit temporal persistence and macroeconomic seasonality across harvest cycles, standard independent and identically distributed (i.i.d.) assumptions severely underestimate standard errors. We structure the 364 calendar-day active trading period into $G = 12$ non-overlapping temporal clusters. We explicitly distinguish two complementary statistical estimands:
\begin{enumerate}
    \item \textit{Pooled Scenario-Level Estimand} ($N = 5,000$): Evaluates dispatch-level net realization differences across all 5,000 pooled scenarios, reported with 2,000-draw paired bootstrap 95\% confidence intervals centered on the pooled mean difference.
    \item \textit{Cluster-Robust Estimand} ($G = 12, \text{df} = 11$): Macro-averages differences across the 12 temporal blocks using the Cameron-Gelbach-Miller cluster-robust sandwich variance estimator:
    \begin{equation}
    V_{\text{cluster}} = \frac{G}{G-1} \frac{N-1}{N-k} (X'X)^{-1} \left( \sum_{g=1}^G X_g' u_g u_g' X_g \right) (X'X)^{-1}
    \end{equation}
    reported with cluster-based 95\% confidence intervals ($\text{df} = 11$), cluster $t$-statistics, and Benjamini-Hochberg False Discovery Rate (FDR) adjustments at $\alpha = 0.05$.
\end{enumerate}

\subsection{Computational Performance \& Complexity}
AgriLink exhibits $O(M)$ time complexity and $O(M)$ space footprint, where $M$ is the number of candidate mandis. Benchmarked on an Intel Core i7-11800H CPU @ 2.30 GHz with 16 GB RAM across 10,000 repeated executions over 18 candidate mandis (including distance calculations, lookback slicing, price aggregation, and candidate ranking), mean execution latency was \textbf{1.14 ms} (median: 1.08 ms, 95th percentile: 1.32 ms). The measured execution latency on the tested Intel Core i7 platform indicates that the decision engine has low computational overhead. Mobile-device performance on low-power rural smartphones remains to be validated in future field deployments.

\section{Empirical Dataset and Regional Context}
\subsection{Data Sources \& Cohort Classification}
We classify the empirical assets into five distinct categories: 1) AGMARKNET administrative records (2,304 validated records across 18 mandis); 2) Synthetic smallholder dispatch scenarios (5,000 dispatches across five seeds); 3) Primary commercial transporter quotes (25 quotes across 4 vehicle tiers); 4) Prospective observational farmer cases (60 smallholder tracking records); 5) FPO organizational pilot adopters (7 FPOs, 32 collective shipments). We emphasize that the 5,000 simulation dispatches represent synthetic scenarios modeled on validated market distributions, completely separate from the 60 human farmer observational tracking records.

\subsection{Preprocessing, Attrition Pipeline \& Evaluation Calendar}
The underlying AGMARKNET administrative repository spans August 15, 2024 to August 14, 2025 (365 calendar days). However, in India, August 15 is Independence Day—a statutory national public holiday during which all regulated APMC wholesale assembly yards are closed for physical auctions. Consequently, the empirical forward evaluation spans the active 364-day marketing calendar (\textbf{August 16, 2024 to August 14, 2025}).

From 5,409 retrieved AGMARKNET records~\cite{agmarknet2024}, state filtering retained 3,982 Andhra Pradesh records; commodity filtering retained 3,118 records across five target crops; daily market-day mean aggregation resolved multi-variety arrivals into 2,741 unique market-days; enforcing the $k \ge 10$ active auction days threshold removed 437 sporadic observations, yielding exactly \textbf{2,304 validated transaction records} across \textbf{18 APMC mandis}.

\section{Simulation Framework and Baseline Suite}
\subsection{Synthetic Farmer Scenario Generation}
We parameterized 1,000 synthetic farmer dispatch scenarios per seed across five independent random seeds (Seeds 42, 101, 2024, 777, 999), generating a pooled cohort of $N = 5,000$ scenarios. Origins were sampled with radial distance $r \sim U(5.0, 35.0)$~km around mandi centroids with polar angle $\theta \sim U(0, 2\pi)$. Parcel sizes were drawn uniformly from $Q \in [1,000, 3,500]$~kg (pooled mean: 2,999.3 kg), reflecting single LCV loads. Eligibility requires $\ge 2$ active candidate mandis within $[T-30, T-1]$. All five seeds use the identical underlying validated AGMARKNET mandi price matrix across different spatial farm coordinates, so they represent spatial-seed stability tests rather than independent empirical replications.

\subsection{Evaluated Baseline Strategy Suite}
We evaluate eight distinct baseline strategies (S1 to S8), AgriLink (Proposed), and an Ex-Post Oracle:
\begin{itemize}
    \item \textbf{S1 (Nearest Market):} $\arg\min_m D_r(L_f, L_m)$. Standard default heuristic.
    \item \textbf{S2 (Naive Highest Price):} $\arg\max_m \bar{P}_{m,c}(T)$. Headline price-seeking heuristic.
    \item \textbf{S3 (Distance-Capped Nearest):} $\arg\min D_r$ subject to $D_r \le 40$~km, with fallback S1. Farm origins are sampled within $r \in [5, 35]$~km geodesic distance. Under road circuity $C_f = 1.28$, road distance to the nearest mandi ranges from 0.87~km to 44.87~km (mean: 24.37~km). In 4,493 scenarios (89.86\%), the nearest mandi has $D_r \le 40$~km and is selected directly. In the remaining 507 scenarios (10.14\%), where the nearest mandi's road distance lies between 40.0~km and 44.87~km, no market satisfies the 40~km threshold, triggering the strategy's explicit fallback rule (\textit{if no candidate is within 40~km, dispatch to the nearest available market}). Consequently, across all 5,000 scenarios, $S_3$ selects the exact same market as $S_1$, producing identical net realizations.
    \item \textbf{S4 (Fixed Cost Hurdle):} Switches from nearest only if $\bar{P}_m \ge \bar{P}_{\text{near}} + \text{₹}2.50/\text{kg}$; otherwise stays at nearest.
    \item \textbf{S5 (Long-Run Historical Mean):} Expanding historical mean computed from start of dataset (Aug 15, 2024) to $T-1$ minus freight. Missing days handled by averaging over all active trading days.
    \item \textbf{S6 (7-Day Moving Average):} Retrospective 7-day mean minus freight.
    \item \textbf{S7 (Robust Worst-Case / Minimax):} Classic Wald maximin decision rule selecting the market that maximizes the conservative minimum price observed over lookback minus freight: $\arg\max_m \{ Q \cdot \min_t P_{m,c}(t) - c_{\text{rate}} \cdot Q \cdot D_r \}$.
    \item \textbf{S8 (Top-3 Restricted Portfolio):} Ranks mandis by $D_r$, selects 3 geographically closest, and chooses the one maximizing expected net realization.
    \item \textbf{AgriLink (Proposed):} 30-day lookback mean minus circuity-corrected freight (Eq. 9).
    \item \textbf{Ex-Post Oracle:} Ex-post optimal market using forward price hindsight $P_{m,c}(T+H)$.
\end{itemize}

\subsection{Forward-Price Evaluation and Imputation}
Forward evaluations are conducted on observed forward prices at primary horizon $H = 7$ days (weekly planning horizon) and categorized into three explicit strata based on forward price availability:
\begin{itemize}
    \item \textbf{Stratum 1 (Zero Imputation):} Directly observed auction trade on exact date $T+H$ ($N = 3,135$ scenarios, 62.70\%).
    \item \textbf{Stratum 2 (Adjacent Window Trade, 1--3d):} Inactive on $T+H$, price drawn from adjacent 1-to-3-day active trading window ($N = 1,550$ scenarios, 31.00\%).
    \item \textbf{Stratum 3 (Last-Price Fallback):} Defaulted to latest pre-dispatch price prior to day $T$ ($N = 315$ scenarios, 6.30\%).
\end{itemize}

\subsection{Statistical Testing Procedure and Temporal Blocks}
The 364 calendar-day active trading study period (August 16, 2024 to August 14, 2025; excluding the August 15 holiday closure) is partitioned into 12 non-overlapping temporal blocks: Blocks 1 to 11 span exactly 30 calendar days each (Block 1: Aug 16--Sep 14; Block 2: Sep 15--Oct 14; Block 3: Oct 15--Nov 13; Block 4: Nov 14--Dec 13; Block 5: Dec 14--Jan 12; Block 6: Jan 13--Feb 11; Block 7: Feb 12--Mar 13; Block 8: Mar 14--Apr 12; Block 9: Apr 13--May 12; Block 10: May 13--Jun 11; Block 11: Jun 12--Jul 11); Block 12 spans the remaining 34 calendar days (Jul 12--Aug 14, 2025) to complete the annual marketing year without truncation ($11 \times 30 + 34 = 364$). All 5,000 synthetic dispatches are mapped to their respective block $g \in \{1, \dots, 12\}$ based on decision date $T$. For each strategy $k$, the cluster-level mean net realization in block $g$ is $\bar{\Pi}_{k, g} = \frac{1}{|\mathcal{S}_g|}\sum_{i \in \mathcal{S}_g}\Pi_k(i)$. The paired cluster-level difference for block $g$ is $d_{g, k} = \bar{\Pi}_{\text{AgriLink}, g} - \bar{\Pi}_{k, g}$. The cluster-robust estimand reports the unweighted macro-mean $\bar{d}_k = \frac{1}{12}\sum_{g=1}^{12} d_{g, k}$, sample standard error $\text{SE}(\bar{d}_k) = \sqrt{\frac{1}{11\times 12}\sum_{g=1}^{12}(d_{g,k}-\bar{d}_k)^2}$, paired cluster $t$-statistic $t_k = \bar{d}_k / \text{SE}(\bar{d}_k)$ with $\text{df} = 11$, and cluster 95\% confidence interval $[\bar{d}_k \pm t_{0.975, 11}\cdot\text{SE}(\bar{d}_k)]$ where $t_{0.975, 11} = 2.201$. Paired 2,000-draw bootstrap confidence intervals evaluate the pooled scenario-level estimand ($N = 5,000$). Multiple testing across strategies is controlled using Benjamini-Hochberg False Discovery Rate (FDR) adjustments at $\alpha = 0.05$.

\section{Empirical Results and Comparative Analysis}
\subsection{Overall Performance Across Cohorts}
Across the pooled 5,000 synthetic dispatches, AgriLink achieves a mean net realization of \textbf{₹153,868.54}, delivering a statistically significant gain of \textbf{+₹13,263.18 (+9.43\%, pooled 95\% CI: [+₹12,634.80, +₹13,887.82]; cluster mean difference +₹10,751.27, cluster $t = 5.62, p < 0.001$, cluster 95\% CI: [+₹6,540.67, +₹14,961.87])} over Nearest Market (S1) and \textbf{+₹1,573.23 (+1.03\%, pooled 95\% CI: [+₹1,380.30, +₹1,752.81]; cluster mean difference +₹1,694.93, $t = 5.30, p < 0.001$, cluster 95\% CI: [+₹991.05, +₹2,398.81])} over Naive Highest Price (S2), confirming \textbf{Hypothesis 1 ($H_1$)}. AgriLink captures \textbf{95.56\%} of the Ex-Post Oracle arbitrage spread with an average absolute regret of only ₹615.68 (0.40\%).

\subsection{Performance by Forward Imputation Stratum}
Stratifying by forward trade availability demonstrates that AgriLink's outperformance is not driven by imputation artifacts: Stratum 1 achieves +6.82\% over S1 and +1.36\% over S2; Stratum 2 achieves +12.13\% over S1 and +0.95\% over S2; Stratum 3 achieves +14.92\% over S1 and +1.33\% over S2.

\subsection{Scenario Outcomes and Negative Arbitrage}
Naive Highest Price (S2) yields lower net realization than Nearest Market (S1) in \textbf{18.94\% of dispatches (947 of 5,000 cases)}, suffering an average net loss of \textbf{-₹8,054.97} per failed dispatch due to excessive transit mileage, proving \textbf{Hypothesis 2 ($H_2$)}. In contrast, AgriLink yields lower net realization than Nearest Market in only \textbf{3.68\% of cases (184 instances)}.

Across all 5,000 dispatches, Nearest Market (S1) incurs a mean transit distance of 24.37~km, AgriLink incurs 113.41~km, and Naive Highest Price (S2) incurs 204.51~km. Relative to Naive Highest Price, AgriLink saves an average of \textbf{91.11 km of one-way transit mileage} across all 5,000 dispatches. In divergent scenarios ($N = 1,222$, 24.44\%) where AgriLink rejects the naive peak-price market, AgriLink saves an average of \textbf{372.78 km of one-way transit mileage} relative to S2 (AgriLink mean haul: 70.22~km vs. S2 mean haul: 443.00~km) and delivers an average net surplus of \textbf{+₹6,437.10 per dispatch}. Relative to Nearest Market, AgriLink incurs an additional 89.04~km of haulage to capture profitable spatial arbitrage opportunities.

\subsection{Commodity Heterogeneity}
Performance varies substantially across commodity categories: Perishable Mango achieves the highest relative gain (+38.07\% over S1) due to high geographic price dispersion. Low-value Maize exhibits modest gains (+7.93\% over S1) because transport expenditure represents a large fraction of gross value. Groundnut achieves +10.39\%, Cotton achieves +2.50\%, and Arhar achieves +7.86\%.

\subsection{Failure-Case Analysis and Boundaries}
AgriLink underperformed Nearest Market in 184 of 5,000 scenarios (3.68\%). A post-hoc audit reveals three primary failure drivers: 1) \textit{Transient Destination Gluts (54.3\% of losses)}: In 100 cases, a distant market reported high historical prices over $[T-30, T-1]$, but experienced an unexpected arrival surge on date $T+H$, depressing clearing prices after transit costs were incurred; 2) \textit{Low Value-Density Vulnerability in Maize (28.8\% of losses)}: Low clearing prices (₹21.65/kg) amplify sensitivity to freight forecasting errors; 3) \textit{Circuity Distortion on Secondary Rural Roads (16.9\% of losses)}: Local road circuity exceeded the regional $C_f = 1.28$ mean, causing freight charges to exceed model expectations.

\section{Robustness, Sensitivity, and Model Ablations}
\subsection{Out-of-Sample Independent Freight Scoring}
To eliminate concerns of circularity, all 5,000 scenarios were scored under one baseline model (Model A) and four alternative, independent freight models. Under stochastic shocks (Model D) and backhaul penalties (Model E), AgriLink's advantage over Naive Highest Price expands from +1.03\% to \textbf{+2.88\% (+₹4,210.71)} and \textbf{+2.22\% (+₹3,310.76)} respectively, confirming \textbf{Hypothesis 3 ($H_3$)}. Naive price chasing routes dispatches across distant district borders, incurring severe freight volatility and empty return penalties.

\subsection{Multi-Seed Simulation Stability}
Spatial stability across five independent seeds ($N = 1,000$ per seed) shows relative gains over Nearest Market ranging from +9.07\% to +9.85\%, and gains over Naive Highest Price ranging from +0.84\% to +1.33\%, with pooled mean net realization of ₹153,868.54.

\subsection{Break-Even Tariffs and Switching Dynamics}
In the primary benchmark seed (Seed 42, $N = 1,000$), AgriLink recommends switching away from the nearest market in exactly \textbf{529 scenarios (52.90\%)}. Across all five pooled spatial seeds ($N = 5,000$), exactly \textbf{2,658 dispatches (53.16\%)} switch markets (Seed 42: 529 [52.9\%], Seed 101: 527 [52.7\%], Seed 2024: 542 [54.2\%], Seed 777: 538 [53.8\%], Seed 999: 522 [52.2\%]).

Figure 6 presents the break-even tariff distribution specifically for the $N = 529$ switched cohort in Seed 42. In this cohort, 510 scenarios (96.41\%) yield positive break-even tariffs with a median of $C^* = \text{₹}0.0720/\text{kg}\cdot\text{km}$ (7.2$\times$ the baseline commercial rate of ₹0.010/kg$\cdot$km), while 19 scenarios (3.59\%) face negative break-even tariffs due to adverse ex-post price reversals. Furthermore, 77.1\% of recommendations survive a 100\% freight tariff increase (₹0.020/kg$\cdot$km), and 68.2\% survive a 200\% increase (₹0.030/kg$\cdot$km).

\subsection{Seasonal Out-of-Sample Stability and Ablations}
Because AgriLink is a deterministic, rule-based optimization engine rather than a trained parametric model, no parameters are fitted or trained on historical subsets. To evaluate seasonal out-of-sample stability, we partitioned the 364-day active trading period into H1 (Kharif harvest: Aug 16, 2024 to Feb 14, 2025; Seed 42 $N = 277$; Pooled 5-Seed $N = 1,431$ [28.62\%]) and H2 (Rabi harvest: Feb 15, 2025 to Aug 14, 2025; Seed 42 $N = 723$; Pooled 5-Seed $N = 3,569$ [71.38\%]). AgriLink delivers +6.40\% over S1 in H1 and +11.07\% over S1 in H2, confirming that spatial arbitrage gains persist across seasonal transitions without degradation.

\section{Primary Field Survey and Case Studies}
\subsection{Primary Commercial Freight Survey}
To ground freight parameters, we surveyed 25 commercial transporters across Kurnool, Kadapa, and Guntur across four vehicle categories: Tata Ace (9 quotes, ₹20.93/t-km), Bolero Maxi Truck (9 quotes, ₹11.46/t-km), Eicher 14ft (5 quotes, ₹9.12/t-km), and Eicher 17ft (2 quotes, ₹7.75/t-km).

\textit{Methodological Justification for Baseline Tariff and Contractual Rate Delineation:} The quote-count pooled mean across all 25 quotes is ₹14.11/t-km, while the unweighted macro-average across the four vehicle categories is ₹12.32/t-km. Mini-trucks charge higher rates (₹20.93/t-km) due to fixed base flag-drops (₹750--₹1,000) over short trips. However, inter-mandi spatial switching in AgriLink occurs over 50--160~km with 1.0--3.5 tonne parcels, exclusively served by Light Commercial Vehicles (Bolero Maxi Truck and Eicher 14ft). The 14 intermediate LCV quotes yield a quote-count weighted average of \textbf{₹10.63/tonne-km} (and an unweighted class average of ₹10.29/tonne-km).

The primary survey identified two distinct transport pricing regimes: individual ad-hoc spot market chartering (where intermediate LCVs charge a quote-count weighted average of ₹10.63/t-km) versus structured seasonal volume charter agreements negotiated between registered Farmer Producer Organizations (FPOs) or aggregators and local transporter unions. In field interviews with FPO management committees in Kurnool and Kadapa, FPOs documented negotiated seasonal volume charter rates operating at \textbf{₹10.00/t-km} ($c_{\text{rate}} = \text{₹}0.010/\text{kg}\cdot\text{km}$), secured by committing multi-trip dispatch volume during peak harvest flushes. We therefore adopt ₹10.00/t-km as our calibrated FPO contract baseline. 

Importantly, the empirical quote benchmark of ₹9.12/t-km for the 4.0-tonne Eicher 14ft vehicle (Table 9) supports the commercial feasibility of the negotiated ₹10.00/t-km rate by demonstrating that dedicated multi-tonne LCV haulage achieves sub-₹10 tariffs on regional corridors. We distinguish these figures: ₹10.00/t-km is the conservative macro-baseline adopted for general multi-seed simulation dispatches (which encompass mixed vehicle consignments and shorter hauls), whereas ₹9.12/t-km is the vehicle-specific empirical rate utilized in dedicated truckload pooling. For unorganized individual smallholders relying solely on spot markets, the relevant benchmark is ₹10.63/t-km. Crucially, sensitivity analysis demonstrates that even when evaluated under unorganized spot tariffs (₹10.63/t-km) or short-haul slab tariffs (₹15.00/t-km), AgriLink maintains robust and expanding net realization gains (+9.12\% and +7.84\% over S1; +1.28\% and +2.15\% over S2) because naive price chasing incurs excessive deadhead transit mileage.

\textit{Vehicle Payload vs Parcel Size:} Smallholder dispatches in the synthetic simulation are parameterized with parcel volumes $Q \sim U(1,000, 3,500)$~kg to represent single-farm harvest deliveries. The Eicher 14ft vehicle has a registered gross payload capacity of 4.0 tonnes (4,000 kg), providing sufficient capacity to accommodate a full 3.5-tonne single consignment with a 12.5\% safety cushion, or to aggregate multiple smaller partial-truckload parcels (e.g., two 1.5--2.0 tonne deliveries from neighbouring smallholders or FPOs) without exceeding axle weight limits.

\subsection{Farmer Participant Selection \& Observational Study Framing}
A prospective cohort of 60 smallholder farmers was tracked across Kurnool, Kadapa, and Guntur during the Kharif 2024 and Rabi 2024--25 harvest seasons. Participants were recruited through local agricultural extension contacts. Because participants were not randomly assigned, this tracking represents an \textbf{observational study} demonstrating real-world associations rather than a randomized controlled trial:
\begin{itemize}
    \item \textbf{Compliant Adopters ($n = 46$):} Farmers who followed AgriLink recommendations achieved an observed net gain of \textbf{+₹14,180.20 (+9.93\%)} over local APMC counterfactuals.
    \item \textbf{Non-Compliant Controls ($n = 14$):} Farmers who defaulted to local yards due to subjective buyer ties incurred an observed opportunity loss of \textbf{-₹13,850.50 (-8.87\%)}.
\end{itemize}

\textit{Evaluation Protocol:} The local APMC counterfactual price was recorded from actual auction modal prices at the farmer's nearest local APMC on the identical sale date for the same commodity and commercial quality grade. The 9.93\% net realization gain represents the percentage difference:
\begin{equation}
\text{Gain}_{\%} = \frac{\Pi_{\text{net, recommended}} - \Pi_{\text{net, local counterfactual}}}{\Pi_{\text{net, local counterfactual}}} \times 100\%
\end{equation}
where net revenues account for gross sale value minus audited transport freight and actual APMC deductions. Transport disbursements and APMC deductions (cess, commission, hamali) were audited from physical market receipts (Form H/J slips) for both cohorts.

\subsection{FPO Logistics Pooling Analysis}
Seven Farmer Producer Organizations (FPOs) participated in a 90-day pilot coordinating 32 collective pooled shipments across 118 tonnes hauled:
\begin{itemize}
    \item \textbf{Individual Baseline Freight:} Individual smallholders chartering small LCVs (Tata Ace, 1.5-tonne payload) at the empirical rate of ₹20.93/t-km incurred an average transport cost of \textbf{₹1,835.56 per tonne} across an average haul of 87.7~km ($87.7\text{ km} \times \text{₹}20.93/\text{t-km} = \text{₹}1,835.56/\text{t}$).
    \item \textbf{Pooled Full-Truckload Freight:} FPOs aggregated parcels into full 4.0-tonne medium LCV loads (Eicher 14ft; average load: 3.69 tonnes per dispatch across 32 dispatches, $\frac{118\text{ t}}{32\text{ dispatches}} = 3.6875\text{ t}$). The total pooled logistics cost averaged \textbf{₹1,348.82 per tonne}, comprising:
    \begin{enumerate}
        \item Vehicle line-haul freight: $87.7\text{ km} \times \text{₹}9.12/\text{t-km} = \text{₹}799.82/\text{tonne}$.
        \item First-mile farm-to-hub aggregation feeder transit: \textbf{₹285.00 per tonne}.
        \item Village hub loading, bagging, and weighing hamali: \textbf{₹128.50 per tonne}.
        \item FPO dispatch coordination fee: ₹500/trip divided by the 3.69-tonne average load = \textbf{₹135.50 per tonne} (or ₹135.59/t across aggregate 118 tonnes).
    \end{enumerate}
    Total pooled cost: $799.82 + 285.00 + 128.50 + 135.50 = \text{₹}1,348.82/\text{tonne}$.
    \item \textbf{Observed Reduction:} Pooling achieved a \textbf{26.52\% freight cost reduction} per tonne: $\Delta_{\text{freight}} = (1835.56 - 1348.82) / 1835.56 = 26.52\%$, saving \textbf{₹486.74 per tonne}.
    \item \textbf{Contractual Rate Basis:} While survey quotes for Eicher 14ft vehicles were collected over 115--160~km corridors, FPOs negotiated structured seasonal charter agreements with local transporter unions based on the ₹9.12/t-km benchmark for the 87.7~km multi-mandi haul.
    \item \textbf{Net Economic Surplus:} Across 118 tonnes hauled, direct transport savings (₹486.74/t) combined with destination price realization premiums (+₹3,060.71/t) generated an aggregate net surplus of \textbf{₹418,600} ($118 \times \text{₹}3,547.45$), averaging \textbf{₹13,081 per dispatch} or \textbf{₹3,547.45 per tonne} of commercialized harvest.
\end{itemize}

\subsection{Research Ethics and Informed Consent}
No formal institutional review board (IRB) or institutional ethics committee (IEC) protocol approval or formal exemption was sought or granted for this study. The research conformed to institutional academic conduct guidelines for non-interventional field research at the School of Computer Science and Engineering (SCOPE), VIT-AP University. The field inquiry was strictly observational, involving only the documentation of routine public agricultural wholesale auction prices, market receipts (Form H/J slips), and commercial freight quotes. No clinical, physiological, psychological, or dietary interventions were performed, and no biological specimens or sensitive personally identifiable data were collected. All participating smallholder farmers and commercial transport operators provided voluntary verbal informed consent following a briefing on the research objectives, and participated without financial compensation. All survey records and market receipts were completely anonymized at the point of data entry.

\section{Discussion and Practical Implications}
\subsection{Digital Public Infrastructure (DPI) Integration}
AgriLink's lightweight algorithmic formulation ($O(M)$ time, 1.14 ms latency) enables direct deployment as an edge decision engine within national platforms like e-NAM and mobile extension portals (e.g., Kisan Suvidha). By embedding GPS-based road distance calculations and vehicle freight lookup tables, public portals can transform passive price broadcasting into active logistics-aware dispatch guidance.

\subsection{FPO Truckload Logistics Aggregation}
Individual smallholders producing sub-truckload harvests (1,000--1,500 kg) face high per-unit freight charges. When deployed by Farmer Producer Organizations (FPOs), AgriLink coordinates collective harvesting schedules, pooling parcels into full 4.0-tonne LCV loads and capturing long-distance arbitrage premiums otherwise accessible only to large commercial traders.

\section{Limitations and Threats to Validity}
We identify five methodological boundaries: 1) \textit{Internal Validity}: Relies on backward-looking 30-day price expectations; sudden demand shocks on auction day cannot be fully anticipated; 2) \textit{External Validity}: Empirical results reflect Andhra Pradesh APMC networks; replication across northern wheat-mustard corridors is required; 3) \textit{Construct Validity}: Assumes commercial LCV availability; in remote interior villages, transport shortages may require spot premiums; 4) \textit{Statistical Conclusion Validity}: Evaluated on synthetic farmer origins drawn uniformly around mandi hubs; clustered inference accounts for temporal correlation; 5) \textit{General Equilibrium Feedback}: Simultaneous adoption by large farmer cohorts would alter destination arrivals and depress clearing prices.

\section{Conclusion and Future Work}
This study demonstrates that digital price dissemination alone is insufficient to optimize farm-gate realization without logistics integration. We distinguish three core takeaways: 1) \textit{Simulation Evidence}: Across 5,000 synthetic dispatches in Andhra Pradesh, AgriLink achieves a mean net realization of \textbf{₹153,868.54}, generating statistically significant gains over Nearest Market (+9.43\%) and Naive Highest Price (+1.03\%), while eliminating 91.11~km of unnecessary transit mileage relative to the Naive Highest Price heuristic (S2) across all dispatches and 372.78~km in divergent scenarios (while incurring an incremental 89.04~km of haulage relative to Nearest Market to capture profitable spatial price arbitrage); 2) \textit{Observational Field Evidence}: Primary tracking indicates compliant farmers achieved +9.93\% higher returns, and FPO truckload pooling yielded a 26.52\% freight reduction; 3) \textit{Future Deployment Validation}: Full causal confirmation will require randomized controlled field trials across multi-state cereal belts, dynamic cooperative fleet dispatching, and real-time cold-chain tracking.

\section*{Acknowledgment}
The authors express their deepest gratitude to their project guide, Prof. Pechetti Mounika, School of Computer Science and Engineering (SCOPE), VIT-AP University, Amaravati, for invaluable guidance, technical review, and academic mentorship throughout this research project. The authors also gratefully acknowledge the School of Computer Science and Engineering (SCOPE) and VIT-AP University for providing computational facilities, administrative support, and research infrastructure.

\begin{thebibliography}{15}
\bibitem{goyal2010information} A.~Goyal, ``Information, direct access to farmers, and rural market performance in central India,'' \emph{Amer. Econ. J. Appl. Econ.}, vol.~2, no.~3, pp. 22--45, Jul. 2010, doi: 10.1257/app.2.3.22.
\bibitem{fafchamps2012impact} M.~Fafchamps and B.~Minten, ``Impact of SMS-based agricultural information on Indian farmers,'' \emph{World Bank Econ. Rev.}, vol.~26, no.~3, pp. 383--414, Sep. 2012, doi: 10.1093/wber/lhr056.
\bibitem{mitra2018asymmetric} S.~Mitra, D.~Mookherjee, M.~Torero, and S.~Visaria, ``Asymmetric information and middleman margins: An experiment with Indian potato farmers,'' \emph{Rev. Econ. Statist.}, vol. 100, no.~1, pp. 1--13, Mar. 2018, doi: 10.1162/REST\_a\_00699.
\bibitem{chatterjee2023market} S.~Chatterjee, ``Market power and spatial competition in rural India,'' \emph{Quart. J. Econ.}, vol. 138, no.~3, pp. 1649--1711, Aug. 2023, doi: 10.1093/qje/qjad004.
\bibitem{takayama1971spatial} T.~Takayama and G.~G. Judge, \emph{Spatial and Temporal Price and Allocation Models}. Amsterdam: North-Holland, 1971.
\bibitem{badiane1998spatial} O.~Badiane and G.~E. Shively, ``Spatial integration, transport costs, and the response of local prices to policy changes in Ghana,'' \emph{J. Develop. Econ.}, vol.~56, no.~2, pp. 411--431, Aug. 1998, doi: 10.1016/S0304-3878(98)00072-8.
\bibitem{sekhar2021agricultural} C.~S.~C. Sekhar, D.~Roy, and Y.~Bhatt, ``Agricultural market reforms and e-NAM: Progress, impediments, and emerging empirical trends,'' \emph{Econ. Polit. Wkly.}, vol.~56, no.~44, pp. 38--47, Oct. 2021.
\bibitem{agricensus2018} Directorate of Economics and Statistics, ``All India report on agriculture census 2015--16,'' Ministry of Agriculture and Farmers Welfare, Govt. of India, New Delhi, India, 2018.
\bibitem{agmarknet2024} Directorate of Marketing and Inspection, ``AGMARKNET portal database on wholesale mandi arrivals and prices,'' Ministry of Agriculture, Govt. of India, 2024. [Online]. Available: https://agmarknet.gov.in
\bibitem{ballou2002selected} R.~H. Ballou, H.~Rahardja, and N.~Sakai, ``Selected country circuity factors for road travel distance estimation,'' \emph{Transp. Res. Part A: Policy Pract.}, vol.~36, no.~9, pp. 843--848, Nov. 2002, doi: 10.1016/S0965-8564(01)00044-1.
\bibitem{kaufman2012leakage} S.~Kaufman, S.~Rosset, C.~Perlich, and O.~Stitelman, ``Leakage in data mining: Formulation, detection, and avoidance,'' \emph{ACM Trans. Knowl. Discov. Data}, vol.~6, no.~4, Art. no.~15, pp. 1--21, Dec. 2012, doi: 10.1145/2382577.2382579.
\bibitem{boscoe2012nationwide} F.~P. Boscoe, K.~A. Henry, and M.~S. Zdeb, ``A nationwide comparison of driving distance versus straight-line distance to hospitals,'' \emph{Prof. Geogr.}, vol.~64, no.~2, pp. 188--196, May 2012, doi: 10.1080/00330124.2011.583586.
\bibitem{negi2018farmers} D.~S. Negi, P.~S. Birthal, D.~Roy, and M.~T.~Khan, ``Farmers' market choices: Do formal channels provide better returns?,'' \emph{Agric. Econ. Res. Rev.}, vol.~31, no. Conf, pp. 85--97, 2018, doi: 10.5958/0974-0279.2018.00025.5.
\bibitem{fackler2001spatial} P.~L. Fackler and B.~K. Goodwin, ``Spatial price analysis,'' \emph{Handb. Agric. Econ.}, vol.~1, Part B, pp. 971--1024, 2001, doi: 10.1016/S1574-0072(01)10025-3.
\bibitem{casaburi2013rural} L.~Casaburi, R.~Glennerster, and T.~Suri, ``Rural roads and local market integration,'' \emph{Amer. Econ. Rev. Pap. Proc.}, vol. 103, no.~3, pp. 262--265, May 2013, doi: 10.1257/aer.103.3.262.
\end{thebibliography}

\begin{IEEEbiography}[{\includegraphics[width=1in,height=1.25in,clip,keepaspectratio]{author1.jpg}}]{Pechetti Mounika}
is an Assistant Professor in the School of Computer Science and Engineering (SCOPE), VIT-AP University, Amaravati, Andhra Pradesh, India. Her research interests encompass machine learning, intelligent optimization, agricultural informatics, and computational decision systems. She served as the faculty project guide and research supervisor.
\end{IEEEbiography}

\begin{IEEEbiography}[{\includegraphics[width=1in,height=1.25in,clip,keepaspectratio]{author2.jpg}}]{P. Kiran Kumar}
(Reg. No. 23BCE9610, Student Member, IEEE) is an undergraduate student in Computer Science and Engineering at the School of Computer Science and Engineering (SCOPE), VIT-AP University, Amaravati, Andhra Pradesh, India. His research interests encompass computational decision support systems, agricultural transport optimization, spatial network modeling, and applied data analytics.
\end{IEEEbiography}

\begin{IEEEbiography}[{\includegraphics[width=1in,height=1.25in,clip,keepaspectratio]{author3.jpg}}]{B. Bhargav}
(Reg. No. 23BCE9623) is an undergraduate student in Computer Science and Engineering at the School of Computer Science and Engineering (SCOPE), VIT-AP University, Amaravati, Andhra Pradesh, India. His research interests include distributed data architectures, rural supply chain modeling, and algorithm optimization.
\end{IEEEbiography}

\begin{IEEEbiography}[{\includegraphics[width=1in,height=1.25in,clip,keepaspectratio]{author4.jpg}}]{S. RavindraNath}
(Reg. No. 23BCE9516) is an undergraduate student in Computer Science and Engineering at the School of Computer Science and Engineering (SCOPE), VIT-AP University, Amaravati, Andhra Pradesh, India. His research interests include spatial data mining, cloud computing architectures, and agricultural informatics.
\end{IEEEbiography}

\begin{IEEEbiography}[{\includegraphics[width=1in,height=1.25in,clip,keepaspectratio]{author5.jpg}}]{S. Kranthi Kowshik}
(Reg. No. 23BCE9611) is an undergraduate student in Computer Science and Engineering at the School of Computer Science and Engineering (SCOPE), VIT-AP University, Amaravati, Andhra Pradesh, India. His research interests include applied optimization algorithms, mobile decision systems, and regional market analytics.
\end{IEEEbiography}

\end{document}
